\documentclass[10pt,oneside]{book}
%%%% Page Info + Commands %%%%%{

%packages
\usepackage{pgfplots}
\usepackage{mathrsfs}
\pgfplotsset{compat=1.18}
\usepackage{amsmath}
\usepackage{amsfonts}
\usepackage{charter}
\usepackage{amsthm,letltxmacro}
\usepackage{amssymb}
\usepackage{fancyhdr}
\usepackage{float}
\usepackage[most]{tcolorbox}
\usepackage{enumerate}
\usepackage{enumitem}
\usepackage[dvipsnames]{xcolor}
\usepackage{changepage}
\usepackage{mdframed}
\usepackage{graphicx}

% SMALL PAGE COMMAND
\newcommand{\smallpage}{  \setlength \oddsidemargin{.25in}
            \setlength \textwidth{5.5in}
            \setlength \topmargin{0in}
            \setlength \textheight{9in}}


% Commands for sets of numbers
\newcommand{\Z}{\mathbb{Z}}\newcommand{\R}{\mathbb{R}}\newcommand{\N}{\mathbb{N}}\newcommand{\Q}{\mathbb{Q}}\newcommand{\C}{\mathbb{C}}
\newcommand{\real}[1]{\text{Re}\left(#1\right)}
\newcommand{\im}[1]{\text{Im}\left(#1\right)}

% gordie spacing and boxing commands
\newcommand{\gspc}{\vspace{0.13cm}} % 0.5 space
\newcommand{\gline}[0]{\gspc\\} % 1.5 space
\newcommand{\gbline}{\gspc\gspc\\} % double space
\newcommand{\gbox}[1]{\fbox{\parbox{\linewidth}{#1}}} % multiline box command
\newcommand{\gmod}[1]{\,(\text{mod}\,#1)}


\newcommand{\gimage}[2]{
\begin{figure}[H]
    \centering
    \includegraphics[width=#2]{#1}
\end{figure}
}

\newcommand{\gimageR}[3]{
\begin{figure}[H]
    \centering
    \includegraphics[width=#2,angle=#3]{#1}
\end{figure}
}

\newcommand{\centerit}[1]{{\begin{center} #1 \end{center}}}
\newcommand{\ul}[1]{\underline{#1}}

\newcommand{\prt}[2]{\frac{\partial#1}{\partial#2}}

%%%%%%%% command for graphics %%%%%%%%%%%%%

%}

% COLOR BOX
\newmdenv[backgroundcolor=gray!8,
  linewidth=0pt, innerleftmargin=0pt, innerrightmargin=0pt,
  skipabove=\baselineskip, skipbelow=\baselineskip
]{coloredadjust}

% PROOF
\LetLtxMacro\oldproof\proof
\let\endoldproof\endproof
\renewenvironment{proof}[1][\proofname]
  {\vspace{0.2cm}\begin{adjustwidth}{2em}{2em}
   \oldproof[#1]}
  {\endoldproof
\end{adjustwidth}}

% PROBLEM
\newenvironment{problem}[1]
  {\vspace{-0.1cm}\begin{coloredadjust}\begin{adjustwidth}{2em}{2em}
   \textit{Problem. }#1}
  {\endoldproof
\end{adjustwidth}\end{coloredadjust}\gspc}

%% DEFINITION
\newtcbtheorem[]{definition}{Definition}{enhanced,
  colback=blue!2, colframe=blue!50!black, coltitle=blue!70!black,
  fonttitle=\bfseries,
  boxrule=0pt, toprule=2pt, bottomrule=1pt, leftrule=1pt, rightrule=1pt, arc=1pt, outer arc=1pt,
  attach title to upper, title={#1},
}{dfn}

%% CRITERION
\newtcbtheorem[]{criterion}{Criterion}{enhanced,
  colback=magenta!3, colframe=magenta!60!black, coltitle=magenta!20!black,
  fonttitle=\bfseries,
  boxrule=4pt, toprule=2pt, bottomrule=1pt, leftrule=1pt, rightrule=1pt, arc=1pt, outer arc=1pt,
  attach title to upper, title={#1},
}{ctn}

%% COROLLARY
\newtcbtheorem[]{corollary}{Corollary}{enhanced,
  colback=orange!3, colframe=orange!80!black, coltitle=orange!50!black,
  fonttitle=\bfseries,
  boxrule=4pt, toprule=1pt, bottomrule=1pt, leftrule=1pt, rightrule=1pt, arc=5pt, outer arc=5pt,
  attach title to upper, title={#1},
}{ctn}

%% THEOREM
\newtcbtheorem[]{theorem}{Theorem}{enhanced,
  colback=green!2, colframe=green!40!black, coltitle=green!20!black,
  fonttitle=\bfseries,
  boxrule=4pt, toprule=3pt, bottomrule=1pt, leftrule=1pt, rightrule=1pt,arc=5pt, outer arc=5pt,
  attach title to upper, title={#1},
}{thm}




% color commands
\newcommand{\cg}[1]{\color{OliveGreen}#1\color{white}}
\newcommand{\cb}[1]{\color{BlueViolet}#1\color{white}}


\newcommand{\cpr}[1]{\left(#1\right)}
%%%%%%%%%%%%% PHYSICS SPECIFIC COMMANDS

\newcommand{\vA}{\vec{\mathbf{A}}}
\newcommand{\vB}{\vec{\mathbf{B}}}
\newcommand{\vC}{\vec{\mathbf{C}}}
\newcommand{\vZ}{\vec{\mathbf{0}}}
\newcommand{\norm}{\vec{\mathbf{n}}}
\newcommand{\pvec}[1]{\left\langle#1\right\rangle}
\newcommand{\ar}{\vec{\mathbf{r}}}
\newcommand{\arhat}{\hat{\mathbf{r}}}
\newcommand{\vv}{\vec{\mathbf{v}}}
\newcommand{\delfull}{\pvec{\prt{}{x}, \prt{}{y},\prt{}{z}}}
\newcommand{\delinput}[3]{\pvec{\prt{}{x}\cpr{#1}, \prt{}{y}\cpr{#2},\prt{}{z}\cpr{#3}}}
\newcommand{\dps}{\displaystyle}
\newcommand{\delmatrix}[3]{\begin{bmatrix}\dps\hat x&\dps\hat y&\dps\hat z\\\dps\prt{}x&\dps\prt{}y&\dps\prt{}z\\\dps#1&\dps#2&\dps#3\end{bmatrix}}

%%%%%%%%%%%%%%%%%%%%%%%%%%%%
%%%%%%%%%%%%%%%%%%%%%%%%%%%%%%
%%%%%%%%%%%%%%%%%%%%%%%%%%%%%
%%%%%%%%%%%%%%%%%%%%%%%%%%%%%
%%%%%%%%%%%%%%%%%%%%%%%%%%%%%%
%%%%%%%%% begin doc

%%%% Page Setup %%%%%%%
\smallpage\sloppy
\pagestyle{fancy}
\lhead{\large{\textbf{HW3 - Theory of Computation}}}
\setlength{\headheight}{10pt}



\begin{document}
\begin{enumerate}
  \item Draw the diagram for an NFA that recognizes:
  \begin{enumerate}
    \item $01^*0$
    \gimage{image.png}{5cm}
    \item $(0\cup 1)1^* (0\cup 1)$
    \gimage{image2.png}{5cm}
    \item $(0\cup 10 \cup 11){(0 \cup 1)}^*$
    \gimage{image3.png}{5cm}
  \end{enumerate}
  \item Give regular expressions describing the following languages. In all cases, the alphabet is $\{0,1\}$
  \begin{enumerate}
    \item $\{ w: \text{1011 is a substring of }w\}$
    \begin{problem}
      I would describe the language with the regular expression:
      \begin{align*}
        \Sigma^*1011\Sigma^*
      \end{align*}
    \end{problem}
    \item $\{w:\text{ the length of }w\text{ is odd}\}$
    \begin{problem}
      All I want to do is ensure that there is $1 + 2k,\;k\in\Z$ elements:
      \begin{align*}
        \Sigma(\Sigma\Sigma)^*
      \end{align*}
    \end{problem}
    \item $\{w:\text{ the number of 1s in }w\text{ is even}\}$
    \begin{problem}
      For this one, all I need to do is ensure that all 1s come in pairs with any number of 0s inbetween them:
      \begin{align*}
        0^*\cpr{10^*1}^*0^*
      \end{align*}
    \end{problem}
  \end{enumerate}
  \item For any string $w = w_1,..., w_n$, the \textbf{reverse} of $w$, written as $w^R$, is the string $w$ in reverse order. For example, the reverse of 0011 is 1100, and the reverse of 0101 is 1010. For any language $A$, let $A^R =\{w^R:w\in A\}$. Show that if $A$ is regular, so is $A^R$. 
  \begin{problem}
    Because $A$ is a regular language, let $A$ be represented by the regular expression $M$. Reverse the ordering of all \textbf{symbols}, treating each part of the expression in parenthesis $()$ as a single symbol. Then, for all elements in parenthesis, repeat this process.\gline{}
    This expression is $A^R$. Thus, because $A^R$ can be represented by a regular expression, $A^R$ is a regular language if and only if $A$ is a regular langage. \gline{}
    \textit{Note: }A \textbf{symbol} in this case is a single character or union of single characters. Any multicharacter expression must be reversed, even in a union. For example, take the expression:
    \begin{align*}
      \cpr{01(0\cup 10)^*\Sigma\Sigma^*01}^R = 10\Sigma^*\Sigma(0\cup 01)10
    \end{align*}
  \end{problem}
  \item Write ``real world'' regular expressions to match the following strings. There are various websites that you can use
  \begin{enumerate}
    \item Dates in the format MM/DD/YY or MM/DD
    \begin{problem}
      Here is the representative regular expression:
        \begin{verbatim}\d{1,2}\/\d{1,2}\/(\d{2}|\d{4})\end{verbatim}
    \end{problem}

    \item Prices in the format \$D.cc. Here, D can be any dollar amount (but should not start with 0). There should always be exactly two digits after the decimal, for cents. For example, your regular expression should match \$.99, \$237.17, and \$12818.00, but not \$0.99, \$010.00, \$58., \$218.2, or \$198484
    \begin{problem}
      Here is the solution I came up with:
      \begin{verbatim}\$([1-9]\d*)*\.\d{2}\end{verbatim}
    \end{problem}
    \item Any ten-digit telephone number \textbf{either} in the form $(xxx)xxx-xxxx$\textbf{ or }$xxx-xxx-xxxx$.
    \begin{problem}
      I used the ? operator to make sure the parenthesis are optional:
      \begin{verbatim}\(?\d{3}(\)|\-)\d{3}\-\d{4}\end{verbatim}
    \end{problem}
  \end{enumerate}
\end{enumerate}

\end{document}